% !TeX root = ../main.tex
\section{Device State and Working Set}\label{sec:device-state}

\begin{hypothesis}
  Back-to-back sustained sequential writes should first ride the SLC/DRAM write cache (a burst), then
  settle to a lower steady-state throughput once the cache must flush to QLC in the background. Occasional
  latency spikes from garbage collection should appear in the upper tail even while mean throughput is
  steady.
\end{hypothesis}

\begin{hypothesis}
  (Added before the large-file runs.) (i) A single sequential write far larger than the SLC cache should
  show a sharp \emph{cliff}: burst throughput near the QD=32 write rate of \cref{sec:block-size}, then a
  drop to the native QLC program rate, several times lower. (ii) After that, random \qty{4}{\kibi\byte}
  writes reach a lower steady state than in the burst, with heavier tails from garbage collection.
  (iii) A short idle period lets the controller drain the cache and recover part of the burst rate.
  (iv) Random reads should slow as the working set grows from \qty{1}{\gibi\byte} to
  \qty{60}{\gibi\byte}, because more logical-to-physical map lookups miss the controller's mapping cache.
  (v) \enquote{Sequential} reads of data written randomly should be much slower than of data written
  sequentially. (vi) Temperature should rise during sustained writes but stay below the
  \qty{86.85}{\celsius} warning threshold reported in sysfs (\code{temp1_max}).
\end{hypothesis}

\begin{knobs}
  \code{sustained} & one O\_DIRECT sequential pass over a fresh \qty{60}{\gibi\byte} file (128\,KiB, QD32;
                     the file is \code{fallocate}d, so every byte is written exactly once). Then
                     \qty{4}{\kibi\byte} random writes at QD32 over the whole file for \qty{180}{\second}
                     or \qty{8}{\gibi\byte}, a \qty{60}{\second} idle gap, and a \qty{15}{\second} /
                     \qty{4}{\gibi\byte} sequential recovery burst. Throughput, p50/p99/p99.9/max latency
                     and NVMe temperature are logged every second \\
  \code{working_set} & \qty{4}{\kibi\byte} random read (QD1, QD32) and random write (QD32, \qty{256}{\mebi\byte}
                     cap) restricted to the first 1, 4, 16 and \qty{60}{\gibi\byte} of the written file \\
  \code{layout}     & on the \qty{2}{\gibi\byte} test file: \qty{128}{\kibi\byte} sequential read (QD1,
                     QD32) and \qty{4}{\kibi\byte} random read, first after a sequential fill, then after
                     a \qty{4}{\gibi\byte} random \qty{4}{\kibi\byte} overwrite (about
                     \qty{86}{\percent} of blocks rewritten in random order) \\
  Volume / safety   & \code{SUSTAIN_GIB=60} (the default 64 would have left less than the required
                     \qty{10}{\gibi\byte} free); the file is deleted afterwards. The filesystem has no online discard, so the
                     freed blocks reach the SSD only at the weekly \code{fstrim} \\
  Drive state       & \code{device_state} started at 10:54, \emph{immediately} after \code{run.sh all}
                     (10:32--10:54, \ensuremath{\sim}\qty{60}{\gibi\byte} written), with no idle period
                     and no TRIM (no root; \code{fstrim} runs weekly) \\
  Script            & \code{SUSTAIN_GIB=60 run.sh device_state} (= \code{sustained} then
                     \code{working_set} on the same file); \code{run.sh layout} (inside \code{all}) \\
\end{knobs}

\subsection{SLC-cache exhaustion and steady state}
\pgfplotsset{susplot/.style={acs, height=0.62\linewidth, legend style={font=\scriptsize}}}
\begin{figure}[tbp]
  \centering
  \IfDataExists{sustained}{%
    \begin{subfigure}[t]{0.49\linewidth}
      \begin{tikzpicture}
        \begin{axis}[susplot, xlabel={GiB written}, ylabel={Write MiB/s (1 s bins)}, ymin=0]
          \addplot[mark=none, blue!70!black] table[x=gib_done, y=mibps, only if={phase}{seq_fill}] {data/sustained.csv};
        \end{axis}
      \end{tikzpicture}
      \caption{Sequential fill: throughput vs.\ volume.}
    \end{subfigure}\hfill
    \begin{subfigure}[t]{0.49\linewidth}
      \begin{tikzpicture}
        \begin{axis}[susplot, xlabel={GiB written}, ylabel={Latency (\si{\micro\second})}, ymode=log,
                     legend pos=north west]
          \addplot[mark=none, black!60] table[x=gib_done, y=lat_p50_us, only if={phase}{seq_fill}] {data/sustained.csv};
          \addlegendentry{p50}
          \addplot[mark=none, red!70!black] table[x=gib_done, y=lat_p999_us, only if={phase}{seq_fill}] {data/sustained.csv};
          \addlegendentry{p99.9}
        \end{axis}
      \end{tikzpicture}
      \caption{Sequential fill: per-second latency.}
    \end{subfigure}

    \begin{subfigure}[t]{0.49\linewidth}
      \begin{tikzpicture}
        \begin{axis}[susplot, xlabel={Time in phase (s)}, ylabel={4 KiB write IOPS}, ymin=0]
          \addplot[mark=none, blue!70!black] table[x=t_s, y=iops, only if={phase}{rand_steady}] {data/sustained.csv};
        \end{axis}
      \end{tikzpicture}
      \caption{Random-write phase (QD32).}
    \end{subfigure}\hfill
    \begin{subfigure}[t]{0.49\linewidth}
      \begin{tikzpicture}
        \begin{axis}[susplot, xlabel={Time (s)}, ylabel={NVMe temperature (\si{\celsius})},
                     restrict y to domain=0:120]
          \addplot[mark=none, orange!80!black] table[x=t_global_s, y=temp_c] {data/sustained.csv};
        \end{axis}
      \end{tikzpicture}
      \caption{Composite temperature, all phases.}
    \end{subfigure}%
  }
  \caption{Sustained-write experiment, one-second samples. (a) The SLC-cache cliff appears as the drop in
  throughput after the burst; (b) the latency tail during the fill; (c) random-write steady state after
  the fill; (d) temperature across the baseline idle, fill, random, idle and recovery phases.}\label{fig:sustained}
\end{figure}

\begin{table}[tbp]
  \centering\small
  \caption{Sustained-write summary (computed by \code{run.sh} into \code{sustained_summary.csv}). Burst is
  the median of the leading seconds that stay above half the first second's rate. The cliff is the first
  second, starting a 5-second window, whose rate falls below half the burst; the table gives the volume
  written before it. Steady is the median of the last \qty{30}{\percent} of the fill. The random phase
  alternates between two states, so the table gives its overall average and the per-second
  quartiles.}\label{tab:sustained}
  \begin{tabular}{@{}lr@{}}
    \toprule
    Burst sequential write (MiB/s) / duration (s)  & \csvval{sustained_summary}{burst_mibps} / \csvval{sustained_summary}{burst_s} \\
    Cliff after (GiB written / s)                  & \csvval[1]{sustained_summary}{cliff_gib} / \csvval{sustained_summary}{cliff_s} \\
    Post-cliff steady sequential write (MiB/s)     & \csvval{sustained_summary}{steady_mibps} \\
    Steady / burst                                 & \csvval[3]{sustained_summary}{steady_over_burst} \\
    Fill: GiB, seconds, average MiB/s              & \csvval[1]{sustained_summary}{fill_gib}, \csvval{sustained_summary}{fill_s}, \csvval{sustained_summary}{fill_avg_mibps} \\
    Random 4 KiB write, QD32: average IOPS         & \csvval{sustained_summary}{rand_steady_iops} \\
    \quad per-second IOPS quartiles; median per-second p99 (\si{\micro\second})
                                                   & \csvval{sustained_summary}{rand_iops_q25}--\csvval{sustained_summary}{rand_iops_q75}; \csvval{sustained_summary}{rand_steady_p99_us} \\
    Recovery after \qty{60}{\second} idle: first / last \qty{5}{\second} (MiB/s) & \csvval{sustained_summary}{recovery_mibps} / \csvval{sustained_summary}{recovery_late_mibps} \\
    NVMe temperature min--max (\si{\celsius})      & \csvval[1]{sustained_summary}{temp_min_c}--\csvval[1]{sustained_summary}{temp_max_c} \\
    \bottomrule
  \end{tabular}
\end{table}

\paragraph{Analysis.} The fill began at the SLC ingest rate,
\csvval{sustained_summary}{burst_mibps}\,MiB/s, for only \csvval{sustained_summary}{burst_s}\,s. After
\csvval[1]{sustained_summary}{cliff_gib}\,GiB it fell to about \qty{62}{\mebi\byte\per\second} and stayed
there for the remaining \qty{57}{\gibi\byte}: the median of the last \qty{30}{\percent} is
\csvval{sustained_summary}{steady_mibps}\,MiB/s, and the 20-minute average is
\csvval{sustained_summary}{fill_avg_mibps}\,MiB/s. A burst rate \num{18}$\times$ the steady rate is far
more than the \enquote{several times} the hypothesis expected. The post-cliff rate is what this QLC drive
sustains when it must program QLC directly while also folding cached data. It has two sub-states. Most
60-second windows sit at \qtyrange{60}{64}{\mebi\byte\per\second}. Eight of the 21 drop to
\qtyrange{28}{46}{\mebi\byte\per\second}, with per-second p99 latencies of
\qtyrange{0.3}{0.5}{\second} at QD32. One isolated second at \qty{97}{\second} wrote \qty{1.03}{\gibi\byte}
at full SLC speed, presumably into a block of cache the controller had just freed.

\textbf{What this run can and cannot say about the cache size.} Only about \qty{3}{\gibi\byte} of SLC
was available when the fill started. That is a property of the run order, not of the drive.
\code{device_state} began immediately after \code{run.sh all} had written
\ensuremath{\sim}\qty{60}{\gibi\byte} in 22 minutes, with no idle time and no TRIM, so the cache was
already nearly full. The fresh-drive burst-then-cliff transition, and hence the full dynamic SLC
capacity, was therefore \emph{not} captured, and we report no cache size. What the rest of the study
shows is the burst \emph{rate}, consistently: \ensuremath{\sim}\qty{1.06}{\gibi\byte\per\second} in the
sequential-write sweep at QD=32 (\cref{sec:block-size}), and 128\,KiB QD=1 sequential writes at
\csvq{fio_check}{experiment=blkbench,mode=write,pattern=seq}{iops_median} IOPS
(\ensuremath{\sim}\qty{1.0}{\gibi\byte\per\second}) in \code{fio_check}, both from different experiments
earlier in \code{all}. The block-size sweep also caught one exhaustion episode, after about
\qty{10}{\gibi\byte} of writes (\cref{sec:block-size}). Measuring the true cache size needs a fill
started on a drive that has idled and been trimmed.

The random-write phase averaged \csvval{sustained_summary}{rand_steady_iops} IOPS
(\ensuremath{\sim}\qty{38}{\mebi\byte\per\second}). This is about \num{22}$\times$ below the burst-state
random-write rate at the same QD (\csvq{queue_depth}{qd=32,mode=write}{iops_median} IOPS,
\cref{sec:queue-depth}), and it alternated between two states. Per-second IOPS quartiles are
\csvval{sustained_summary}{rand_iops_q25} and \csvval{sustained_summary}{rand_iops_q75}: stretches near
\num{16}k IOPS, where writes still land in freshly freed SLC, alternate with stretches near \num{3}k
IOPS, where foreground writes wait for garbage collection. The median per-second p99 is
\qty{32}{\milli\second}. After a \qty{60}{\second} idle there was \emph{no} recovery. The recovery burst
ran at \csvval{sustained_summary}{recovery_mibps}\,MiB/s for its first seconds and fell to
\csvval{sustained_summary}{recovery_late_mibps}\,MiB/s by the end. With tens of gigabytes waiting to be
folded into QLC, a minute of idle frees essentially no SLC, so hypothesis (iii) fails at this time
scale.

Two caveats. First, the file covers \qty{60}{\gibi\byte} of a \qty{512}{\giga\byte} drive, and the
filesystem does not trim on delete. This is steady state for this region at this fill level, not the
SNIA full-drive steady state, which would need the raw device (ruled out by the data-safety rule).
Second, the drive started this run in a used state, as described above.

\paragraph{Temperature.} The composite temperature ranged from \csvval[1]{sustained_summary}{temp_min_c}
to \csvval[1]{sustained_summary}{temp_max_c}\,\si{\celsius} during this experiment, and from
\qtyrange{33.9}{61.9}{\celsius} over all runs (the peak was during the QD=32 sequential-read sweep).
Both are far below the \qty{86.85}{\celsius} warning and \qty{89.85}{\celsius} critical thresholds
reported by the drive. The throughput drop happened within seconds, at \qty{50.9}{\celsius}, and did
not track temperature afterwards (flat at \qty{55.9}{\celsius} while throughput alternated), so it is
the cache, not thermal throttling.

\subsection{Working set (LBA range)}
\begin{figure}[tbp]
  \centering
  \IfDataExists{working_set}{%
    \begin{subfigure}[t]{0.49\linewidth}
      \begin{tikzpicture}
        \begin{axis}[acs, height=0.85\linewidth, xmode=log, log basis x=2, xtick={1,4,16,60}, xticklabels={1,4,16,60},
                     xlabel={Working set (GiB)}, ylabel={QD=1 read mean latency (\si{\micro\second})}]
          \addplot+[iqr bars] table[x=span_gib, y=lat_mean_us_median, iqr of=lat_mean_us, only if={mode}{read}, only if={qd}{1}] {data/working_set.csv};
        \end{axis}
      \end{tikzpicture}
      \caption{QD=1 random read latency.}
    \end{subfigure}\hfill
    \begin{subfigure}[t]{0.49\linewidth}
      \begin{tikzpicture}
        \begin{axis}[acs, height=0.85\linewidth, xmode=log, log basis x=2, xtick={1,4,16,60}, xticklabels={1,4,16,60},
                     xlabel={Working set (GiB)}, ylabel={QD=32 IOPS}, ymin=0, legend pos=south west]
          \addplot+[iqr bars] table[x=span_gib, y=iops_median, iqr of=iops, only if={mode}{read}, only if={qd}{32}] {data/working_set.csv};
          \addlegendentry{rand read}
          \addplot+[dashed, iqr bars] table[x=span_gib, y=iops_median, iqr of=iops, only if={mode}{write}, only if={qd}{32}] {data/working_set.csv};
          \addlegendentry{rand write}
        \end{axis}
      \end{tikzpicture}
      \caption{QD=32 random read and write IOPS.}
    \end{subfigure}%
  }
  \caption{\qty{4}{\kibi\byte} random I/O restricted to the first 1--\qty{60}{\gibi\byte} of the written
  file (medians, IQR bars). The runs followed the sustained write, so the drive was still folding in the background..}\label{fig:working-set}
\end{figure}

QD=1 random-read latency goes from
\csvq[1]{working_set}{span_gib=1,mode=read,qd=1}{lat_mean_us_median}\,\si{\micro\second} at
\qty{1}{\gibi\byte} to \csvq[1]{working_set}{span_gib=4,mode=read,qd=1}{lat_mean_us_median} (4),
\csvq[1]{working_set}{span_gib=16,mode=read,qd=1}{lat_mean_us_median} (16) and
\csvq[1]{working_set}{span_gib=60,mode=read,qd=1}{lat_mean_us_median}\,\si{\micro\second} at
\qty{60}{\gibi\byte}, a rise of about \qty{37}{\percent}. QD=32 read IOPS fall from
\csvq{working_set}{span_gib=1,mode=read,qd=32}{iops_median} to
\csvq{working_set}{span_gib=60,mode=read,qd=32}{iops_median} (\qty{-25}{\percent}). The direction fits
logical-to-physical map lookups missing the controller's mapping cache more often as the span grows: at
4 bytes per \qty{4}{\kibi\byte} page, \qty{60}{\gibi\byte} needs \qty{60}{\mebi\byte} of map. Most of the
increase happens between 1 and \qty{4}{\gibi\byte}, so the cached portion of the map is small. We could
not tell whether the drive uses a host memory buffer (no allocation message in the kernel log).

Two confounds limit this conclusion. First, every working-set point ran right after the sustained
experiment, with tens of gigabytes still being folded. Even the \qty{1}{\gibi\byte} span is far slower
than the \csvq[1]{layout}{experiment=layout_seq_filled,bs=4096}{lat_mean_us_median}\,\si{\micro\second}
QD=1 random read on the \qty{2}{\gibi\byte} test file earlier, so absolute values describe a busy drive.
Second, the spans ran in increasing order, so the trend is also confounded with elapsed time. Random
writes (\csvq{working_set}{span_gib=1,mode=write,qd=32}{iops_median},
\csvq{working_set}{span_gib=4,mode=write,qd=32}{iops_median},
\csvq{working_set}{span_gib=16,mode=write,qd=32}{iops_median} and
\csvq{working_set}{span_gib=60,mode=write,qd=32}{iops_median} IOPS for spans of 1, 4, 16 and
\qty{60}{\gibi\byte}) show no span trend. They jump between the same two garbage-collection states as
the random-write phase above, which supports the expectation that write cost does not depend on logical
address once the cache is full.

\subsection{Data layout: why preconditioning order matters}\label{sec:layout}
\begin{table}[tbp]
  \centering\small
  \caption{The same reads on the same \qty{2}{\gibi\byte} file after a sequential fill and after a random
  \qty{4}{\kibi\byte} overwrite (medians of 5 repetitions).}\label{tab:layout}
  \begin{tabular}{@{}lrr@{}}
    \toprule
    Read & after sequential fill & after random overwrite \\
    \midrule
    \qty{128}{\kibi\byte} seq, QD=1 (MiB/s)  & \csvq{layout}{experiment=layout_seq_filled,bs=131072,qd=1}{mibps_median}
                                              & \csvq{layout}{experiment=layout_rand_overwritten,bs=131072,qd=1}{mibps_median} \\
    \qty{128}{\kibi\byte} seq, QD=32 (MiB/s) & \csvq{layout}{experiment=layout_seq_filled,bs=131072,qd=32}{mibps_median}
                                              & \csvq{layout}{experiment=layout_rand_overwritten,bs=131072,qd=32}{mibps_median} \\
    \qty{4}{\kibi\byte} rand, QD=1 (IOPS)    & \csvq{layout}{experiment=layout_seq_filled,bs=4096}{iops_median}
                                              & \csvq{layout}{experiment=layout_rand_overwritten,bs=4096}{iops_median} \\
    \bottomrule
  \end{tabular}
\end{table}

An early pilot run, since discarded, exposed this effect by accident: a \qty{128}{\kibi\byte} QD=1
sequential read run right after a random-write test was several times slower than the block-size sweep.
An SSD places data in the order it was \emph{written}, not by logical address. \Cref{tab:layout} tests
this directly, and the hypothesis is confirmed strongly. After the random overwrite, QD=1 sequential
reads drop from \csvq{layout}{experiment=layout_seq_filled,bs=131072,qd=1}{mibps_median} to
\csvq{layout}{experiment=layout_rand_overwritten,bs=131072,qd=1}{mibps_median}\,MiB/s ($6.4\times$).
\qty{4}{\kibi\byte} random reads are unaffected
(\csvq{layout}{experiment=layout_seq_filled,bs=4096}{iops_median} vs.\
\csvq{layout}{experiment=layout_rand_overwritten,bs=4096}{iops_median} IOPS), since a single page costs
the same wherever it sits.

Our sub-prediction that QD=32 would hide the gap was \emph{wrong}. At QD=32 the fragmented file gives
\csvq{layout}{experiment=layout_rand_overwritten,bs=131072,qd=32}{mibps_median}\,MiB/s against
\csvq{layout}{experiment=layout_seq_filled,bs=131072,qd=32}{mibps_median}, a $9.2\times$ gap, and no
better than at QD=1. Each fragmented \qty{128}{\kibi\byte} request becomes up to 32 scattered page
reads. The drive completes \ensuremath{\sim}\num{1830} such requests per second (about \num{58}k page
reads/s) at either depth, well below the \num{130}k random \qty{4}{\kibi\byte} reads/s it sustains at
QD=32. We hypothesise that the controller serializes the fragments of one command, or that read-ahead
actively wastes bandwidth on the wrong pages, but we did not test this.

This is the \enquote{insufficient preconditioning} pitfall in a subtler form. A sequential benchmark is
meaningful only if the data's write history is controlled. That is why \code{run.sh} rewrites the test
file sequentially before every sequential-read group. One group missed this in the run reported here:
in \code{all}, \code{fio_check} followed \code{layout} without a rewrite, so its sequential points
measured the fragmented layout (\cref{sec:fio-check}). The script has since been fixed.
